\documentclass[12pt]{article}
\usepackage{amsmath} 
\usepackage{esint} 
\usepackage{amsfonts} 
\usepackage{amssymb} 
\usepackage{geometry} 
\usepackage{mathrsfs}

\usepackage[style=numeric,backend=biber]{biblatex}
\addbibresource{references.bib}


\newcommand{\R}{\mathbb{R}}
\newcommand{\atlas}{\mathscr{A}}
\newcommand{\p}{\partial}
\newcommand{\at}[2]{\left. #1 \right|_{#2}}
\newcommand{\bv}[1]{\mathbf{#1}}
\newcommand{\uv}[1]{\hat{#1}}

\newcommand{\ihat}{\hat{\imath}}
\newcommand{\jhat}{\hat{\jmath}}
\newcommand{\khat}{\hat{k}}



\begin{document}

\title{Exterior Calculus}
\author{Daniel Pesic}
\maketitle
\newpage
\tableofcontents
\newpage
\section{Introduction from Vector Calculus}
\subsection{Vector Calculus Motive}
Vector Calculus is a powerful yet limited tool. It allows us to extend the concepts from single-variable Calculus to higher dimensional spaces, more complex functions, and different varieties of functions. However, its methods become inadequate when space is globally non-Euclidean or when the functions we study require a more general framework; when we ask: 
\begin{itemize}
    \item How can we analyze a function with a locally Euclidean input space rather than a globally Euclidean one? 
     \item In Vector Calculus, curl is limited to vector fields in $\R^3$. How can we find the vorticity of a higher dimensional vector field?
    \item How can we analyze functions intrinsically without relying on coordinates?
\end{itemize}

The primary goal of Exterior Calculus is to provide an abstract framework with a specialized algebra to analyze complex geometric systems. It must not replace Vector Calculus, but build upon it to help us understand what lies beyond. % something else to reference exterior algebra
\subsection{Limitations of Integral Theorems}
Consider the three integral theorems of Vector Calculus. These theorems are viewed as the capstone of the field, giving a geometric relation between the ideas of flux, vorticity, boundaries, and integrals. 
\\\\ Green's Theorem
\[
  \oint_{\partial D} \bigl(P\,dx + Q\,dy\bigr)
  \;=\;
  \iint_{D} \left(\frac{\partial Q}{\partial x} - \frac{\partial P}{\partial y}\right) dA
\]
where $D \subset \mathbb{R}^2$ is a simply connected region with positively oriented
boundary $\partial D$, and $P, Q$ are components of a $C^1$ vector field
with $P, Q : \mathbb{R}^2 \to \mathbb{R}$.
\\\\ Stokes' Theorem
\[
  \oint_{\partial S} \bv{F} \cdot d\bv{r}
  \;=\;
  \iint_{S} \bigl(\vec{\nabla} \times \bv{F}\bigr) \cdot \uv{n} \, dS
\]
where $S \subset \mathbb{R}^3$ is an oriented piecewise-smooth surface with unit normal
$\uv{n}$, and $\partial S$ is its boundary curve with orientation induced by $\uv{n}$.
$\bv{F} = \langle P, Q, R \rangle : \mathbb{R}^3 \to \mathbb{R}^3$ is a $C^1$
vector field, and $d\bv{r}$ is the tangent vector element along $\partial S$.
\\\\ Gauss' Divergence Theorem 

\[
  \oiint_{\partial V} \bigl(\bv{F} \cdot \uv{n}\bigr) \, dS
  \;=\;
  \iiint_{V} \bigl(\vec{\nabla} \cdot \bv{F}\bigr) \, dV
\]
where $V \subset \mathbb{R}^3$ is a bounded solid region, $\partial V$ is its
closed boundary surface with outward-pointing unit normal $\uv{n}$,
and $\bv{F} = \langle P, Q, R \rangle : \mathbb{R}^3 \to \mathbb{R}^3$ is a $C^1$
vector field.
\bigskip

From the Vector Calculus viewpoint, these theorems each have specific use cases and seem quite different compared to each other; they are only applicable in sub-regions of $\R^3$ with specific kinds of functions. However, one might notice a pattern shared between them: the integral of a function over a region's boundary equals the integral of function's derivative over a region. This can even be seen in the ordinary Fundamental Theorem of Calculus:

\[
f(b)-f(a) = \int_{[a,b]}{\frac{df}{dx}dx} \iff\int_{\p[a,b]}{f} =\int_{[a,b]}{\frac{df}{dx}dx}
\]

This pattern is not simply coincidental, as we will soon discover how these four theorems emerge from one overarching theorem at the top of Exterior Calculus: the Generalized Stokes' Theorem. 

\subsection{Roadmap}
Exterior Calculus provides a unified framework that expands the concepts of Vector Calculus through extension, generalization, and unification.
\subsubsection{Extension}
Vector Calculus handles derivatives and integrals in $\R^n$. Exterior Calculus extends these concepts to arbitrary dimensions spaces known as manifolds that are locally look like $\R^n$ but may have non-Euclidean behavior globally.
\subsubsection{Generalization}
The derivative in Exterior Calculus, the exterior derivative, generalizes gradient, curl, and divergence into one operator. Similarly, integration over manifolds generalizes line, surface, and volume integrals in a dimension-independent manner. 
\subsubsection{Unification}
Exterior Calculus unifies Green's, Stokes', and Gauss's theorems, along with the Fundamental Theorem of Calculus into one abstract theorem: the Generalized Stokes Theorem.

These functions of Exterior Calculus are precisely what make it such a powerful framework. Throughout this paper, we will begin at the level of Vector Calculus and build up through these extensions, generalizations, and unifications to show the full theory of Exterior Calculus along with its geometric and physical applications.


\section{Manifolds}

A manifold is an $n$-dimensional structure that is locally Euclidean, i.e., locally homeomorpic to $\R^n$. This means that a manifold can be split into a number of regions that can each be mapped to a region of $\R^n$ and back in a way that preserves its local topological structure. For example, each hemisphere of the sphere $S^2$ can be mapped to $\R^2$ and mapped back while preserving its structure. 

Manifolds generalize the notion of space in Exterior Calculus. The primary difference is that manifolds need not be globally Euclidean, only locally Euclidean. Subsets of $\R^n$ are manifolds as well, but globally non-Euclidean manifolds are the cases which separate their techniques from Vector Calculus, so we will focus on them. For Exterior Calculus to work, both functions and manifolds must be smooth; we now examine what that smoothness condition requires of manifolds.

\subsection{Constructing Smooth Manifolds}

\subsubsection{Topological Manifolds}
A manifold is, at its core, a space that looks locally like ordinary Euclidean space. A topological manifold ensures this resemblance at the level of continuity, while a smooth manifold strengthens this resemblance so that calculus can be performed consistently.

In Exterior Calculus, we require the smooth structure so that differentiation is globally well defined. The full technical construction of smooth manifolds is provided in APPENDIX.

\subsubsection{Coordinate Charts}
The core parts of a manifold $M$ that define its geometric relation to Euclidean space are the coordinate charts. These are precisely what defines the locally Euclidean characteristic of manifolds, which gives the basis for analysis. 

Around every point in $M$, there exists a neighborhood $U$ that can be assigned coordinates $(x^1,x^2, \dots, x^n)$ just like a region of $\R^n$. These coordinates give a way to label points in the neighborhood and map them to $\R^n$ while preserving the geometric structure. The coordinate functions $x^i:U\to \R$ give a method of labeling points in the neighborhood; however, it is an important distinction that coordinates are not directly connected to the geometric structure of the manifold. This will be covered in Section 3. The coordinate chart is defined by $(\phi, U)$ where $\phi$ is the map from $M$ to $\R^n$ via the coordinate functions.

However, coordinate charts alone are not enough to describe the full differentiable structure of a manifold. They do not let us understand the structure of differentiability that lies between the charts.

\subsubsection{Transition Maps}
When two coordinate neighborhoods overlap, there must be a way to translate coordinates from one chart to the other. These coordinate translations are called transition maps. They ensure that the transition between the two charts is fully smooth by creating a map that retains the geometric information from the charts. A transition map between two charts $(\phi, U)$ and $(\psi, V)$ is defined by 
\[
 \psi \circ \phi^{-1}: \phi(U \cap V) \to \psi(U \cap V)
\].
Since both $\phi(U \cap V)$ and $\psi(U \cap V)$ are subsets of $\R^n$, smoothness of the transition map is well defined. If the mapping is smooth, it ensures that derivatives in one chart correspond to derivatives in the other. Therefore, the two charts are smoothly compatible with each other. This allows smoothness to become an intrinsic property of the manifold rather than depending on a chart.

It is an important note that if $U \cap V = \varnothing$, the differentiability condition becomes vacuously true. 

\subsubsection{Atlases}

A smooth atlas is defined as a collection of charts on $M$ that cover the entire manifold. If an atlas is smooth, then the collection $\atlas$ of all charts smoothly compatible with each member of the atlas is both maximal and smooth. Therefore, $\atlas$ is a maximal smooth atlas. The proof for this can be viewed in APPENDIX. The maximal smooth atlas creates a smooth structure on $M$ and describes its topology in relation to $\R^n$.

\subsubsection{Smooth Manifolds}

A smooth manifold $M$ is defined as a set of points combined with a maximal smooth atlas $\atlas$. Since smooth manifolds are ubiquitous in Exterior Calculus, we will simply use the term ``manifold'' to refer to one, with ``atlas'' used to refer to its maximal smooth atlas. With this definition, we can extend the idea of space to include all manifolds beyond $\R^n$. By defining a smooth manifold, smooth functions can be defined beyond $\R^n$ and behave well under analysis. 

With our generalized notion of space, we now work to define principles of calculus. Since manifolds only resemble Euclidean space locally, we must define derivatives in a way that does not depend on a particular coordinate system.

\subsection{Tangent Vectors as Derivations}
In Vector Calculus, the directional derivative finds the derivative of a scalar function on a particular vector. In Exterior Calculus, we can define abstract vectors (that is, vectors not defined in coordinates) that are tangent to a manifold and act as a derivative operator, a derivation.
\subsubsection{Local Coordinate Vector Fields}
In $\R^n$, the global coordinates induce a global coordinate vector field. For example, in $\R^3$, the coordinates $(x,y,z)$ induce $\ihat, \jhat, \khat$.

Since a smooth manifold does not have a global coordinate chart, the coordinate vector fields must be expressed locally. Take a point $p$. The coordinate chart
\[
(U,\varphi), \quad \varphi = (x^1,x^2, \dots, x^n), \quad p\in U
\]
induces the local coordinate vector fields $\frac{\p}{\p x^1}, \frac{\p}{\p x^2},\dots \frac{\p}{\p x^n}$ on $U$. These vectors point in the direction of their corresponding coordinate. This notation of expressing coordinate vectors as partial derivative operators will be explained in Section 2.2.3. % Gabe will review notation
\subsubsection{Tangent Vectors}
A tangent vector on a manifold is a vector attached at a specific point $p$ such that it is tangent to the surface of the manifold. In coordinates, a tangent vector is a linear combination of the local coordinate vector fields. However, we will define a more geometric definition for abstract tangent vectors that does not rely on an arbitrary coordinate chart.

Choose a smooth parametric curve on an embedded manifold $\gamma:I \subseteq \R \to M$ such that $\gamma(0)=p$. A vector $v$ is a tangent vector at $p$ if and only if $\gamma'(0) = v$. An important distinction is that the derivative of a curve on a manifold is only predefined when the manifold can be embedded in $\R^n$. We can choose a more abstract definition using derivations. 

\subsubsection{Derivations}
A derivation is a differential operator that computes directional derivatives along an abstract tangent vector. The tangent vector is not an input to a derivation, but the operator itself. Specifically, a tangent vector is a derivation $v:C^{\infty}{(M)} \to \R$.

For an operator to be a derivation, it must satisfy these two axioms:
\begin{enumerate}
    \item Linearity: $\forall f,g \in C^{\infty}(M)$ and $c,d \in \R, \quad v(cf + dg) = cv(f)+dv(g)$
    \item Leibniz Rule: $v(fg) =fv(g)+gv(f)$
\end{enumerate}
In coordinates, the tangent vector becomes a linear combination of local coordinate vectors:
\[
v^i\frac{\p}{\p x^i}
\]
The numerical value of the derivation applied to a function is:
\[
v^i\at{\frac{\p f}{\p x^i}}{p} 
\]
The use of partial derivative operators as local coordinate vectors gives a natural way to encode this concept. It causes the derivation to act like a directional derivative, similarly to the Euclidean directional derivative:
\[
\bv{v} \cdot \at{\nabla f}{p} = v^i \at{\frac{\p f}{\p x^i}}{p}
\]
Having partial derivative operators as the tangent vector basis lets the tangent vector fundamentally act as a derivation without connecting it to an external operator such as $\nabla$.

Another way to represent the derivation on a function is through the abstract definition with parametric curves:
\[
v(f)=\frac{d}{dt}\at{f(\gamma(t))}{t=0}
\]
This generalization through tangent vectors gives us an abstract way to represent directional derivatives on manifolds and serves as a crucial stepping stone towards higher level concepts.
\subsubsection{More on the Leibniz Rule}
The Leibniz rule $v(fg)=fv(g)+gv(f)$ is a minimal equation that encodes differential properties into a functional. It creates a derivation by satisfying these axioms of differentiation:
\begin{enumerate}
    \item The derivative of a constant must be 0
    \item Derivatives must only extract first order change
\end{enumerate}

When a constant is plugged into a derivation, the Leibniz rule forces:
\[
v(1\cdot1) = 1\cdot v(1)+1\cdot v(1) \implies v(1)=2v(1) \implies v(1)=0
\]
\[
\forall c\in\R, v(c)=v(\sum_{c}{1}) = \sum_{c}{v(1)} = \sum_{c}{0} = 0
\]
Additionally, the Leibniz rule forces derivations to extract solely first order change. Consider the product rule of differentiation for functions in $\R^n$, the specific case of the Leibniz rule in Euclidean space:
\[
\frac{d}{dx^n}\left[fg\right] =f\frac{dg}{dx^n}+g\frac{df}{dx^n}
\]
If both $f$ and $g$ are 0 at a point $p$, the derivative in every direction must be 0. Regardless of higher order change and curvature, the function $fg$ must have no first order change, ensuring that the derivation measures only the instantaneous change of a function at a point. This property is retained when generalizing to manifolds, which enforces the idea how the Leibniz rule forces derivations to satisfy the axioms of differentiation.

\subsection{Tangent Spaces and Bundles}

 While tangent vectors are useful as derivations, we can group them into ordered structures to help us better understand a manifold.

\subsubsection{Tangent Spaces}
The set of all tangent vectors at a point $p$ is called the tangent space $T_pM$. Since, in coordinates, every tangent vector at a point is a linear combination of its local coordinate vectors $\frac{\p}{\p x^1},\frac{\p}{\p x^2}...\frac{\p}{\p x^n} $, the local coordinate vectors form a basis for $T_pM$ with respect to a local chart. Every tangent vector $v$ at $p$ satisfies $v \in T_pM$.

On an $n$-manifold, $T_pM$ is an $n$-dimensional vector space isomorphic to $\R^n$

\subsubsection{Tangent Bundles}
The tangent bundle $TM$ is the set of all tangent vectors on a manifold. It is defined by: \[
TM :=\bigsqcup_{p\in M}{T_pM}
\]
The disjoin union $\sqcup$ is defined to pair each vector with its origin point, functioning as:
\[
\bigsqcup_{p\in M}{T_pM} = \bigcup_{p\in M}{\{p\} \times T_pM}
\]
Therefore, each element of $TM$ takes the form $(p,v)$. This allows us to create a projection map:
\[
\pi:TM \to M \quad(p,v) \mapsto p
\]
Rather than working with each $T_pM$ in isolation, assembling them into $TM$ creates a global structure over $M$. The projection $\pi$ assigns every vector in $TM$ to its base point while $\pi^{-1}(p) \cong T_pM$ recovers the vector fiber over $p$. 

Specifically, $\pi^{-1}$ maps $p$ to the set $\{(p,v) : v \in T_pM\}$ which is naturally isomorphic to $T_pM$, giving a way to obtain the tangent space from a point. Additionally, since both tangent spaces are isomorphic to $\R^n$, the labeling inherited from the disjoint union ensures that for $p \neq q, T_pM \neq T_qM$. Labeling each vector with its base point ensures that global mappings and assignments can exist where each point has a unique tangent space. 

\subsubsection{Vector Fields}
A vector field, that is, a section of the tangent bundle, assigns a tangent vector to each point on a manifold: 
\[
\bv{X}:M\to TM
\]
It must satisfy $\pi \circ \bv{X} = \text{id}_M$, where $\text{id}_M$ is the identity map on the manifold. This ensures that vectors are correctly assigned, as the vector field must satisfy $\bv{X}(p)= v\in T_pM$ for all $p \in M$. This uniqueness allows for the assignments to represent geometric quantities, as a vector $\bv{X}(p)$ will not belong to the tangent space of another point.

\subsection{Submanifolds and Boundaries}
The final parts of manifolds we must mention are submanifolds and boundaries. These are imperative to define future notions of integration.

\subsubsection{Submanifolds}

Given two manifolds $M$ and $N$, $M$ is a submanifold of $N$ if and only if $M \subseteq N$. For example, the sphere $S^2$ is a submanifold of $\R^3$ as $S^2 \subset \R^3$. Submanifolds give us a generalized way to define regions for integration. For example, the region $[0,1] \times [0,1]$ can be expressed as a submanifold of $\R^2$. 

There are also two more structured types of submanifolds: immersions and embeddings. 

An immersion is a smooth map between manifolds $f:M \to N$ such that the image $f(M) \subseteq N$ is locally injective. 
$f$ is locally injective if and only if for every open subset $U$ of $M$, $\at{f}{U}$ is injective. 
Because the image of the immersion $f(M)$ is not globally injective, it may intersect itself, and its topology may not match $M$. For example, $M = (0,2\pi)$ can be immersed in $\R^2$ through $f:M \to \R^2, \quad f(t) \mapsto (\cos(t),\sin(t))$ as a circle missing one point. $0$ and $2\pi$ are distant from each other in $M$, but they become infinitely close in $f(M)$. Therefore, sequences that diverge in $M$ can converge in $f(M)$, so they are not topologically equivalent.

An embedding is an immersion $f:M \to N$ such that $f$ is globally injective and $f(M) \subseteq N$ is homeomorphic to $M$. Due to these properties, $f(M)$ cannot intersect itself and is topologically equivalent to $M$.

APPENDIX covers more on how manifolds can be immersed or embedded in Euclidean space.

\subsubsection{Boundaries}

In simplest terms, the boundary of a smooth $n$-manifold $M$ is a smooth $(n-1)$-manifold $\p M$ embedded as its boundary points. In this case, $M$ is a manifold with boundary while $\p M$ is a manifold without boundary. The notation $\p M$ refers to the boundary of M. In essence, the boundary of a manifold is the set of points on the manifold that are not in the interior.

 Formally, the interior of $M$, $\text{int}(M)$, is defined as the set of points contained in an open subset of $M$. The boundary of $M$ is the set of all points contained within open subsets $p \in U \subset M$ such that:
 \[
U \cap \text{int}(M) \neq \varnothing
 \]
 \[
 U \setminus \text{int}(M) \neq \varnothing
 \]

In essence, the boundary of a manifold is the set of points on the manifold that are not in the interior. 

Boundary manifolds are deeply connected to boundary integrals, as they create the region that is integrated over. By defining boundary manifolds, we generalize boundary integration to manifolds, which is crucial when we discuss integration.

\section{Coordinate Invariance}

Coordinates are used frequently in fields such as Vector Calculus to describe and quantify change. However, there are infinitely many coordinate systems that can be used to describe a structure such as a vector space or manifold. Coordinate invariance dictates that a certain object encodes the same geometric properties regardless of the coordinate system used to describe it. For example, $(1, 1)$ in Cartesian coordinates represents the same geometric vector as $( \sqrt{2}, \frac{\pi}{2} )$ in polar coordinates.

\subsection{Coordinates as Convenience}

Coordinates are incredibly useful for deriving quantitative information from a geometric system. However, they do not constitute the geometry itself. A sphere, such as the Earth, has intrinsic properties such as distance and curvature. Latitude and longitude are one system used to describe these properties, but spherical coordinates, Cartesian coordinates, or any other smooth coordinate chart can be used to describe them. Regardless of which label is chosen to represent a point or a path, the Earth remains unchanged. Therefore, distance and curvature are coordinate invariant, which is precisely what makes them geometric properties.

This does not mean that coordinates are meaningless, however. They are auxiliary tools used to describe a geometric system.

\subsection{Objects Under Coordinate Change}

When coordinates are changed, the numerical components used to describe a geometric object will generally change as well. However, if the object is truly geometric, this change in components must occur in a way that preserves the underlying object itself. This requirement leads to the tensor transformation law.

\subsubsection{Tensor Transformations}

Let $(x^1, \dots, x^n)$ and $(x^{1'}, \dots, x^{n'})$ be two coordinate systems on the same region of a manifold. A vector can be expressed in either coordinate system as
\[
V = V^{i}\frac{\p}{\p x^i} = V^{i'}\frac{\p}{\p x^{i'}}
\]
where $\frac{\p}{\p x^i}$ and $\frac{\p}{\p x^{i'}}$ are local coordinate vector fields. These coordinate vectors are related by the chain rule
\[
\frac{\p}{\p x^j} = \frac{\p x^{i'}}{\p x^j}\frac{\p}{\p x^{i'}}
\]
which gives the Jacobian as the core tool to change coordinates. A conceptual way to consider this is that the tangent basis vectors encode first order behavior as shown in Section 2.2.4. Therefore, the Jacobian, a linear map, contains all of the information needed for the vectors to transform, even if the coordinate change is nonlinear. This transformation ensures that the vector itself remains unchanged, even though its numerical description differs. Objects whose components transform according to this type of rule are called tensors, which include scalars, vectors and more. These objects and their properties will be covered in more detail in Section 4.1.

\subsubsection{Geometric Meaning of the Tensor Transformation Law}

The tensor transformation law provides the precise criterion for distinguishing intrinsic geometric objects from coordinate artifacts. If an object's components transform according to the tensor law, the object exists independently of coordinates and represents an intrinsic geometric quantity. If not, the object depends on the coordinate system and does not represent intrinsic geometry. In this way, tensors form the natural basis of coordinate invariant geometry.

Some objects contain artifacts of the particular coordinate system they are defined in and are referred to as coordinate dependent objects. These are covered more in APPENDIX.

\subsection{The Guiding Principle of Coordinate Invariance}

By prioritizing our analysis on coordinate independent objects, we create a generalized system that allows us to describe the geometry of any given objects without relying on coordinates. While coordinates can be very useful for describing some phenomena, every definition will be given in a coordinate independent manner. This allows us to consider the underlying geometry, which becomes incredibly important as concepts progress. 

\section{Differential Forms}

The central goal of Exterior Calculus is to extend integration to manifolds. In Vector Calculus, integration is defined separately for each type of region: a line integral uses the integrand $\bv{F} \cdot d\bv{r}$, a surface integral uses $\bv{F} \cdot d\bv{S}$, and a volume integral uses $f\,dV$. On a manifold, we need a single, coordinate-independent object that serves as the integrand for regions of any dimension.

A differential form is the generalization of the integrand from Euclidean space to manifolds. A differential $k$-form is an object that can be integrated over $k$-dimensional regions: a $1$-form integrates over curves, a $2$-form over surfaces, and so on. The integrands in Vector Calculus already match this format. Differential forms make the structure explicit and extend it beyond $\R^3$. For differential forms to serve as integrands, they must encode orientation as well as magnitude. Reversing the direction of traversal must flip the sign, just as swapping the bounds of a definite integral does. This requirement will shape the algebraic structure developed in this chapter.

\subsection{Multilinear Maps and Tensors} 

Before we define differential forms as integrands, we need a language that can precisely describe how quantities depend on multiple directions at once. Since a $k$-form operates in $k$ dimensional tangent space and return a scalar, we need a multilinear map from the tangent space to $\R$. Tensors, a kind of multilinear map, provide this exact framework. From this, differential forms will emerge naturally as a special kind of tensor designed specifically for integration.

\subsubsection{Multilinear Maps}

Recall the definition for a linear map between vector spaces $T: V \to W$

\[
T(a\bv{u} + b\bv{v}) = aT(\bv{v})+bT(\bv{u}) \quad a,b \in \R, \; \bv{u},\bv{v} \in V.
\]

A multilinear map is a map from multiple vector spaces $V_n$ 
\[
T: V_1 \times \dots \times V_n \to W
\]

that is linear with respect to each variable. That is, for each variable $u$ in $T$, if all other variables are held constant, $T$ is linear in $u$. For example, the dot product in $\R^n$ is multilinear (specifically bilinear), as
\[
(a\bv{u} + b\bv{v}) \cdot \bv{w} = a(\bv{u} \cdot \bv{w}) + b(\bv{v} \cdot \bv{w}) \text{ and } \bv{u} \cdot (a\bv{v} + b\bv{w}) = a(\bv{u} \cdot \bv{v}) + b(\bv{u} \cdot \bv{w}).
\]

Multilinear maps are a crucial concept as they allow us to define tensors, which is the fundamental concept that differential forms are built upon.

\subsubsection{Covectors}

Before discussing tensors, we must mention covectors. Given a vector space $V$, its dual space $V^*$ is the set of all linear mappings
\[
\alpha: V \to \R.
\]
Each element $\alpha$ is a covector. Therefore, a covector can be viewed as a multivariable, linear function that maps vectors to real numbers.

For example, the map $\alpha: \R^2 \to \R$ defined by $\alpha(\bv{v}) = 3v_1 - v_2$ is a covector.

As we will see, a differential 1-form assigns a covector to each point of a manifold. Therefore, a covector is precisely what a 1-form looks like at a single point.

\subsubsection{Tensors} % say what V_1 and Vq and indicies and blah are V_q -> V_p 2 indicies? im not gonna know what this means but basically add more detail to (p,q) be clear what things are
A tensor is a specific type of multilinear map that maps covectors and vectors to a scalar value.
\[
T: V^*_1 \times \dots \times V^*_p \times V_1 \times \dots \times V_q \to \R.
\]

This is specifically referred to as a type $(p,q)$ tensor, where $p$ counts the covector inputs and $q$ counts the vector inputs. A tensor takes both kinds of input because covectors and vectors pair naturally: given $\alpha \in V^*$ and $\bv{v} \in V$, evaluating $\alpha(\bv{v})$ produces a scalar. A $(p,q)$ tensor extends this pairing to $p + q$ inputs in a multilinear way.

For example, a covector linearly maps one vector to a real number. Therefore, it is a $(0,1)$ tensor via
\[
T: V_1 \to \R.
\]
A vector $\bv{v} \in V$ is a $(1,0)$ tensor. This is the least obvious example, since vectors are not typically viewed as functions. However, given a covector $\alpha \in V^*$, evaluating $\alpha(\bv{v}) \in \R$ is a natural operation. Holding $\bv{v}$ constant and changing $\alpha$, the map $\alpha \mapsto \alpha(\bv{v})$ is linear in $\alpha$. Therefore, $\bv{v}$ is a linear functional on $V^*$:
\[
T: V^*_1 \to \R.
\]
The Euclidean inner product, bilinearly maps two vectors to a real number. Therefore, it is a $(0,2)$ tensor
\[
T:V_1 \times V_2 \to \R.
\]
All of these geometric objects have a representation as a multilinear map. This is the abstract form of a tensor. The coordinate representation can be obtained by applying the tensor to basis elements. For example, in $\R^2$, the dot product of vectors $\bv{x}$ and $\bv{y}$ with basis $\{\ihat, \jhat\}$ can be represented by
\[
T(x_1 \ihat + x_2 \jhat,\; y_1 \ihat + y_2 \jhat).
\]
This can be expressed as a matrix
\[
\begin{bmatrix} 
T(\ihat, \ihat), T(\ihat, \jhat) \\
T(\jhat, \ihat), T(\jhat, \jhat)
\end{bmatrix}.
\]
Therefore, the Euclidean inner product can be expressed as
\[
\begin{pmatrix}
x_1 & x_2
\end{pmatrix}
\begin{bmatrix}
1 & 0 \\
0 & 1
\end{bmatrix}
\begin{pmatrix}
y_1 \\
y_2
\end{pmatrix}
\]

These properties of tensors allow them to describe properties in space, which makes them a natural choice to represent a differential form. Specifically, a differential form is a kind of covariant tensor field, meaning that it purely maps tangent vectors to a real number at each point of a manifold.

\subsection{Differential 1-Forms}

With tensors and covectors in hand, we now define the first kind of differential form. A differential 1-form $\omega$ on a manifold $M$ is a smooth assignment of a covector to each point: at every $p \in M$, $\omega_p$ is a linear functional on the tangent space $T_pM$. This is the generalization of a covector to a manifold. Rather than operating on a single vector space, a 1-form operates across the entire tangent bundle by assigning a covector to each tangent space smoothly.

A scalar function $f: M \to \R$ is a 0-form. Since a 0-form integrates over points, it is integrated by simply evaluating the function at the given point, consistent with the $k$-form framework established at the start of this chapter.

The most natural source of 1-forms is differentiation. Given a smooth function $f: M \to \R$, the differential $df$ is a 1-form that generalizes the gradient from Vector Calculus. In $\R^n$, the gradient $\nabla f$ encodes the rate of change of $f$ in every direction; $df$ does the same on a manifold, but as a covector field rather than a vector field. At each point $p$, $df_p$ assigns to each tangent direction $\bv{v}$ the rate of change of $f$ along $\bv{v}$. Applying this to each coordinate function $x^i: M \to \R$ produces the basis differentials $dx^i$.

The symbols $dx^i$ play two roles simultaneously: they are the differentials of the coordinate functions $x^i$, and they are the basis covectors dual to the coordinate vector fields $\frac{\p}{\p x^i}$. These are the same object. In coordinates, any 1-form $\omega$ on $M$ can therefore be expressed as
\[
\omega = a_i \, dx^i
\]
where $a_i: M \to \R$ are smooth functions and the $dx^i$ span the cotangent space $T_p^*M$ at each point.

\subsection{The Wedge Product and $k$-Forms}

A 1-form integrates over curves. To integrate over a surface, we need an object that takes two tangent vectors as input, one for each dimension of the surface. More generally, to integrate over a $k$-dimensional region we need a $k$-form: an object that takes $k$ tangent vectors and returns a scalar. The wedge product is the operation that builds $k$-forms from 1-forms.

\subsubsection{The Wedge Product}

Given two 1-forms $\alpha$ and $\beta$, their wedge product $\alpha \wedge \beta$ is a 2-form defined by
\[
(\alpha \wedge \beta)(\bv{u}, \bv{v}) = \alpha(\bv{u})\beta(\bv{v}) - \alpha(\bv{v})\beta(\bv{u}).
\]
This combines two 1-forms into an object that maps two tangent vectors to a scalar corresponding to the area of the parallelogram spanned by the vectors. A $k$-form is built by wedging $k$ 1-forms together:
\[
\omega = dx^{i_1} \wedge \cdots \wedge dx^{i_k}.
\]
Any $k$-form on a manifold can be expressed as a linear combination of such products.

\subsubsection{Antisymmetry and Orientation}

The wedge product is antisymmetric, meaning that swapping any two factors flips the sign,
\[
\alpha \wedge \beta = -\beta \wedge \alpha.
\]
Integrating over a surface requires choosing an orientation, and reversing that orientation must flip the sign of the result, just as swapping the bounds of a definite integral does. The antisymmetry of the wedge product encodes this property. Under any permutation $\sigma$ of the factors, the sign changes by $\operatorname{sgn}(\sigma)$:
\[
dx^{i_{\sigma(1)}} \wedge \cdots \wedge dx^{i_{\sigma(k)}} = \operatorname{sgn}(\sigma)\, dx^{i_1} \wedge \cdots \wedge dx^{i_k}.
\]
One property that emerges from this is nilpotency: $\alpha \wedge \alpha = 0$. The permutation must change the sign, but it leaves the form unchanged, forcing the value to be zero:
\[
\alpha \wedge \alpha = - \alpha \wedge \alpha \iff \alpha \wedge \alpha = 0.
\]
 Geometrically, this reflects the fact that a parallelogram spanned by two identical vectors has no area.

\subsubsection{Connection to $n$-Volume}

A $k$-form is multilinear and antisymmetric. Since these are precisely the axioms of the determinant, evaluating a $k$-form on $k$ tangent vectors must compute one:
\[
(dx^{i_1} \wedge \cdots \wedge dx^{i_k})(\bv{v}_1, \dots, \bv{v}_k) = \det \begin{pmatrix} dx^{i_1}(\bv{v}_1) & \cdots & dx^{i_1}(\bv{v}_k) \\ \vdots & \ddots & \vdots \\ dx^{i_k}(\bv{v}_1) & \cdots & dx^{i_k}(\bv{v}_k) \end{pmatrix}.
\]
Each entry $dx^{i_a}(\bv{v}_b)$ is the $i_a$-th component of $\bv{v}_b$. Since the determinant measures oriented $k$-volume, a $k$-form is a smooth field of oriented $k$-volume measurements on the manifold.

\section{Pushforwards and Pullbacks}

In Vector Calculus, the Jacobian allows for the change of variables. It measures how a map changes space and transforms the integrand accordingly. To integrate on manifolds, we need the same capability, but generalized for differential forms. This section builds the framework that allows us to move functions and geometric objects across smooth maps. They will be central to defining integration on manifolds and ultimately, the Generalized Stokes' Theorem.

\subsection{Pullback of Scalar Functions}

Consider a smooth map between manifolds $F: M \to N$. This mapping between spaces can also map objects within them. Consider a function $g: N \to \R$. Given that $g$ is a smooth function on $N$, we can produce a smooth function on $M$ by composing with $F$:
\[
F^* g = g \circ F : M \to \R.
\]
The function $F^* g$ is named the pullback of $g$ by $F$. The map $F$ routes each point $p \in M$ to $F(p) \in N$, then $g$ assigns a scalar value. The result is a function that evaluates $g$ from points on $M$.

This operation, the pullback of scalar functions, is the foundation for moving objects across manifolds, providing a generalized framework for integrand transformations and change of variables. 

\subsection{The Pushforward}

Consider a smooth map $F: \R^m \to \R^n$. The Jacobian matrix at a point $p$ gives the best linear approximation of $F$ near $p$. Given a vector $\bv{v}$ at $p$, the Jacobian maps it to a vector at $F(p)$ by
\[
J_F(p)\,\bv{v} =
\begin{pmatrix}
\dfrac{\p F^1}{\p x^1} & \cdots & \dfrac{\p F^1}{\p x^m} \\[6pt]
\vdots & \ddots & \vdots \\[2pt]
\dfrac{\p F^n}{\p x^1} & \cdots & \dfrac{\p F^n}{\p x^m}
\end{pmatrix}
\begin{pmatrix} v^1 \\ \vdots \\ v^m \end{pmatrix}.
\]
Since $J_F(p)$ is a linear map from $\R^m$ to $\R^n$, it carries the vector space structure of vectors at $p$ to vectors at $F(p)$. This describes how $F$ transforms directions across the vector spaces.

However, one might notice the Jacobian matrix is defined entirely in terms of coordinates. Since manifolds do not have global coordinate systems, the Jacobian cannot be directly applied to maps between manifolds. To generalize this idea, we need a coordinate independent version.

Since tangent vectors on a manifold are derivations, the coordinate-independent generalization of the Jacobian must act on derivations. Consider a smooth map $F: M \to N$ where $p \in M$ is sent to $F(p) \in N$. A tangent vector $v \in T_pM$ acts on smooth functions $f: M \to \R$; we want to produce a derivation at $F(p)$ acting on smooth functions $g: N \to \R$. The differential $dF_p: T_pM \to T_{F(p)}N$ accomplishes this by
\[
(dF_p \cdot v)(g) = v(g \circ F)
\]
for any smooth $g: N \to \R$. Since $g \circ F: M \to \R$ is smooth, $v(g \circ F)$ is well-defined, and $dF_p \cdot v$ satisfies linearity and the Leibniz rule, it is a derivation at $F(p)$. This mapping is called a pushforward.
The pushforward is a linear map between the vector spaces $T_pM$ and $T_{F(p)}N$, matching how the Jacobian behaves between $\R^m$ and $\R^n$. Written in local coordinates $(x^1, \dots, x^m)$ on $M$ and $(y^1, \dots, y^n)$ on $N$, the pushforward acts on coordinate basis vectors via
\[
dF_p\!\left(\frac{\p}{\p x^i}\right) = \frac{\p F^j}{\p x^i}\,\frac{\p}{\p y^j},
\]
where $F^j = y^j \circ F$ are the coordinate components of $F$. The matrix $\left(\frac{\p F^j}{\p x^i}\right)$ is exactly the Jacobian of $F$ in these charts. Therefore, $dF_p$ becomes the intrinsic object and the Jacobian is its coordinate representation.

The name ``pushforward'' comes from how the map acts geometrically: $dF_p$ ``pushes'' each tangent vector at $p$ forward through $F$ to a tangent vector at $F(p)$, just as the Jacobian pushes vectors forward through a map between Euclidean spaces. The notation $F_*$ is used to represent the set of all maps $dF_p \;\; \forall p \in M$.

Recall that the term $g \circ F$ is the pullback of $g$ by $F$. Since $F^{*}g$ is a scalar function on $M$, the derivation $v \in T_pM$ may operate on it, allowing for vectors to be applied to functions across manifolds. In this way, the pushforward of a vector is defined in terms of a scalar function's pullback.

\subsection{Pullback of Differential Forms}

A $k$-form on $N$ operates on $k$ tangent vectors at a point of $N$ and returns a scalar value. We cannot always push a form forward along a general smooth map, though, since a map from $T_{F(p)}N$ to $T_pM$ only exists if $F$ is a homeomorphism. However, with the pushforward $dF_p: T_pM \to T_{F(p)}N$ mapping derivations to $N$, we can map each input vector through the pushforward, then operate with $\omega$. For $v_1, \dots, v_k \in T_pM$, define
\[
(F^*\omega)_p(v_1, \dots, v_k) = \omega_{F(p)}(dF_p \cdot v_1, \dots, dF_p \cdot v_k).
\]
This is the pullback of $\omega$ by $F$. It is a $k$-form on $M$, pulled back along $F$ just like a scalar function.

In local coordinates $(x^1, \dots, x^m)$ on $M$ and $(y^1, \dots, y^n)$ on $N$, the pullback of a coordinate 1-form is
\[
F^*(dy^j) = dF^j = \frac{\p F^j}{\p x^i}\,dx^i,
\]
where $F^j = y^j \circ F$ is the $j$-th component of $F$. For a general $k$-form $\omega = f_{i_1 \cdots i_k}\, dy^{i_1} \wedge \cdots \wedge dy^{i_k}$ on $N$, linearity gives
\[
F^*\omega = (f_{i_1 \cdots i_k} \circ F)\, F^*(dy^{i_1}) \wedge \cdots \wedge F^*(dy^{i_k}).
\]
The Jacobian entries $\frac{\p F^j}{\p x^i}$ appear in each factor, and the wedge product assembles them into the $k$-form coefficients on $M$.

The pullback of differential forms allows us to transform them across manifolds. This is particularly significant as it allows us to simplify and evaluate integrals, eventually culminating in the Generalized Stokes' Theorem.

\section{Integration}
\section{Exterior Derivative}
\section{Generalized Stokes' Theorem}

\section{Appendix}

\subsection*{Evaluation of Basis 1-Forms on Tangent Vectors}

Given a basis 1-form $dx^i$ dual to the coordinate vector field $\frac{\p}{\p x^i}$ at a point on a manifold, its evaluation on a coordinate vector $\frac{\p}{\p x^j}$ is given by
\[
dx^i \left( \frac{\p}{\p x^j}\right) = \delta_{ij}
\]
where
\[
\delta_{ij} = \begin{cases}
1 & \text{if } i = j \\
0 & \text{if } i \neq j
\end{cases}
\]
is the Kronecker delta. Therefore, a basis 1-form extracts a single coordinate direction from a tangent vector. For example,
\[
dx^1 \left( \frac{\p}{\p x^1}\right) = 1, \qquad dx^1 \left( \frac{\p}{\p x^2}\right) = 0.
\]
A general tangent vector $\bv{v} = v^i \frac{\p}{\p x^i}$ is mapped by linearity to $dx^j(\bv{v}) = v^j$: the basis 1-form $dx^j$ simply reads off the $j$-th component of $\bv{v}$.

Another important remark about immersions and embeddings is their connection with the Whitney immersion theorem and the strong Whitney embedding theorem. The Whitney immersion theorem states that any smooth $n$-dimensional manifold can be immersed in $\R^{2n-1}$, while the strong Whitney embedding theorem says that any smooth $n$-dimensional manifold can be embedded in $\R^{2n}$. An excellent example of this is the Klein bottle, a 2-manifold. The Klein bottle can exist in $\R^3$ as an immersion, however, it must intersect itself to stay closed. It can only be embedded in $\R^4$ where it stays closed without intersection.


A coordinate dependent object is an artifact from the given coordinate system and does not represent the geometry of the overall space. For example, consider the second partial derivative of a function with respect to a particular coordinate. By the chain rule, the first derivative transforms as

\[
\frac{\p f}{\p x^{i'}}
=
\frac{\p x^k}{\p x^{i'}}
\frac{\p f}{\p x^k}.
\]

Taking another derivative, we obtain

\[
\frac{\p}{\p x^{j'}} \left( \frac{\p f}{\p x^{i'}} \right) = \frac{\p}{\p x^{j'}} \left(
\frac{\p x^k}{\p x^{i'}}
\frac{\p f}{\p x^k} \right).
\]

When evaluated with the product rule, we obtain
\[
\frac{\p^2 f}
{\p x^{i'} \p x^{j'}}
=
\frac{\p x^k}{\p x^{i'}}
\frac{\p x^l}{\p x^{j'}}
\frac{\p^2 f}{\p x^k \p x^l}
+
\frac{\p^2 x^k}
{\p x^{i'} \p x^{j'}}
\frac{\p f}{\p x^k}
\]

This result does not obey the tensor transformation law, and contains a coordinate artifact. This means that second partial derivatives, by themselves, do not define intrinsic geometric objects. The underlying issue is that the partial derivative operator
\[
\frac{\p}{\p x^i}
\]
	
is tied to the particular coordinate system $x^i$. When the coordinates change, these operators change in a way that reflects the coordinate grid rather than the intrinsic geometry. This is why the second partial derivative is not intrinsically geometric, as it combines the function's intrinsic curvature with the curvature of the coordinate grid.



\section {Citations}
\nocite{*}
\printbibliography

\end{document}

